\chapter{Exploration des données (EDA)}
\label{chap:exploration}

Cette exploration a été menée dans le notebook \texttt{notebooks/analyse\_exploratoire\_TOS.ipynb},
exécuté de bout en bout sans erreur (dernière relecture complète : 10/08/2026, ré-exécutée depuis zéro
plutôt que de faire confiance aux résultats d'une exécution antérieure). Elle constitue le socle
technique de la Phase 0.

\section{Problèmes de qualité des données identifiés}

\begin{itemize}
    \item \textbf{Espacement des pixels incohérent} : le patient \texttt{2017-110\_01-0222-V1MR}
        (les 4 séries) et \texttt{2017-110\_01-0223-V1MR} (série FLASH D uniquement) ont un
        \texttt{PixelSpacing} de 0,7812\,mm contre 1,0\,mm pour le reste --- nécessite un
        rééchantillonnage avant entraînement.
    \item \textbf{Erreur de nommage} : \texttt{2017-110\_01-0222-V1MR/RSSFP D} possède un fichier
        \texttt{Segmentation\_CLAV\_K1.seg.seg.nrrd} (double \texttt{.seg}) et il manque entièrement
        \texttt{Segmentation\_costoclavicular\_mask.nrrd}.
    \item \textbf{\texttt{SliceThickness} non fiable} : varie de façon incohérente entre images d'une
        même série dynamique, probablement altéré à l'export --- ne pas utiliser.
    \item \textbf{Incohérence de casse et fautes de frappe} dans les noms de segments
        (\texttt{K1}/\texttt{k1}, \texttt{VSC}/\texttt{vsc}, \texttt{MCS} au lieu de \texttt{MSC}) ---
        nécessite une normalisation (majuscules + table d'alias) avant toute agrégation de
        statistiques.
    \item \textbf{Deux anomalies structurelles réelles} (non corrigibles par renommage) :
        \texttt{2017-110\_01-0223-V1MR RSSFP D} présente un layout 4D à 2 couches (K1 seul en couche 1,
        les 5 autres structures en couche 0 avec PB=5 au lieu de 6 dans ce fichier) ;
        \texttt{2017-110\_01-0222-V1MR FLASH D} stocke PB à la valeur de label 7 au lieu de 6. Le code
        doit toujours résoudre les structures par leur nom dans l'en-tête, jamais par une valeur de
        label codée en dur.
\end{itemize}

\section{Bug critique n°1 : ordre des images inversé dans le labelmap}
\label{sec:bug-ordre-frames}

L'axe des images (\emph{frames}) dans \texttt{Segmentation\_*.seg.nrrd} est stocké dans l'ordre
\textbf{inverse} de l'ordre \texttt{InstanceNumber}/nom de fichier DICOM : pour une série de $N$
images, l'indice de tableau $k$ correspondant à l'image $n$ est $k = N - n$, et non $n - 1$ comme on
pourrait le supposer naïvement.

Confirmé en comparant \texttt{space origin} / \texttt{space directions} du nrrd avec
\texttt{ImagePositionPatient} de chaque image DICOM. Toute association image~$\leftrightarrow$~label
doit recalculer ce mapping par série (fonction \texttt{build\_frame\_to\_k}) plutôt que de supposer
$k = n-1$.

\textbf{Cinq séries n'ont aucune information spatiale exploitable} pour résoudre cet ordre :
\texttt{2017-110\_01-0222-V1MR} (les 4 séries) et \texttt{2017-110\_01-0223-V1MR/FLASH D} ---
\texttt{ImagePositionPatient} et \texttt{SliceLocation} y sont identiques sur les 20 images, contre
des valeurs variables (et donc exploitables) dans les 51 autres séries. \texttt{build\_frame\_to\_k}
signale correctement ces cas comme non résolus. \textbf{Ces 5 séries doivent être exclues de tout
entraînement automatisé} jusqu'à vérification manuelle sous 3D Slicer.

\textbf{Convergence notable (relecture du 10/08/2026)} : ces 5 séries à géométrie non résolue sont
\emph{exactement} les 5 séries à \texttt{PixelSpacing} anormal (\S précédente) --- pas une coïncidence
a priori, probablement un même lot d'acquisition ou d'export distinct des 51 autres séries. Ces 5
séries doivent être traitées comme un seul groupe anormal à part entière, pas comme deux problèmes
indépendants.

\section{Bug critique n°2 : transposition H/W entre nrrd et DICOM}
\label{sec:bug-transposition}

\textbf{Découvert le 09/08/2026, par inspection visuelle de l'utilisatrice.} Au-delà du problème
d'ordre des images (\S précédente), les deux premiers axes (hauteur/largeur) du labelmap nrrd étaient
\textbf{transposés} par rapport au \texttt{pixel\_array} DICOM.

\subsection{Pourquoi ce bug n'a pas été détecté plus tôt}

\begin{itemize}
    \item La vérification géométrique de l'axe des images (\S précédente) ne portait que sur cet
        axe-là, pas sur les axes H/W.
    \item Une simple comparaison de \emph{shape} ne peut jamais détecter ce bug : les deux axes H et W
        valent 320 dans ce jeu de données, donc une transposition ne change pas la forme du tableau.
        C'est un angle mort structurel de toute vérification d'alignement basée uniquement sur la
        forme --- y compris dans le notebook \texttt{exploration.ipynb} (cellule 48), qui ne compare
        que des \emph{shapes}.
    \item Le raisonnement géométrique initial (correspondance des cosinus directeurs
        \texttt{ImageOrientationPatient} vs \texttt{space directions} du nrrd) suggérait à tort
        qu'aucune transposition n'était nécessaire.
\end{itemize}

\subsection{Méthode de validation retenue}

Un premier contrôle visuel zoomé avait semblé ``plausible'' (structures osseuses dans une zone
anatomiquement cohérente), ce qui s'est révélé insuffisant. La correction a finalement été confirmée
de façon \textbf{quantitative} : un score d'alignement contour/gradient (magnitude du gradient de
Sobel de l'image, moyennée le long du contour du masque) a été calculé pour plusieurs orientations
candidates (identité, transposée, symétries). L'orientation correcte obtient un score systématiquement
plus élevé que l'orientation erronée (p.\,ex.\ 327,9 contre 171,4 sur le premier lot de séries testé).

\textbf{Re-vérification indépendante (10/08/2026)} : le même test a été rejoué sur 5 séries
supplémentaires, couvrant 2 patients non testés auparavant et un mélange FLASH/RSSFP. L'orientation
transposée reste systématiquement gagnante, avec un écart de score de 1,5 à 2,4$\times$ selon la série
--- la correction ne dépend donc pas d'un choix ponctuel de séries de test.

\textbf{Correction appliquée} : ajout de \texttt{.transpose(1, 0, 2)} en fin de la fonction
\texttt{load\_canonical\_labelmap} (désormais dans \texttt{src/data\_loading.py}).

\subsection{Leçon méthodologique}

Ni le raisonnement géométrique à partir des en-têtes DICOM/nrrd, ni un contrôle visuel unique, ne sont
suffisants pour valider un alignement image/masque sur ce projet. Toute future vérification
d'alignement doit s'appuyer sur un test quantitatif, indépendant de l'orientation, et prendre au
sérieux tout retour visuel contredisant une conclusion géométrique.

\section{Vérification de la résolution du masque}
\label{sec:resolution-masque}

Contrôle indépendant du bug de transposition ci-dessus : le \texttt{.seg.nrrd} a-t-il la même
résolution que son image DICOM source ? Deux critères vérifiés sur les 56 séries : les dimensions
H$\times$W du tableau nrrd contre \texttt{Rows}/\texttt{Columns} DICOM, et l'espacement physique
(norme des vecteurs \texttt{space directions} du nrrd) contre \texttt{PixelSpacing} DICOM.

\textbf{Résultat : 56/56 séries concordantes sur les deux critères}, y compris le fichier 4D
multi-couches (\texttt{2017-110\_01-0223-V1MR/RSSFP D}) une fois un piège de parsing corrigé --- dans
un nrrd 4D, l'axe non spatial (\emph{layer}) est représenté par une ligne \texttt{[nan, nan, nan]}
dans \texttt{space directions}, pas par un scalaire \texttt{None} ; un filtre naïf sur le type laisse
passer cette ligne par erreur et produit un faux mismatch. Conclusion : pas de bug de résolution
indépendant du bug de transposition --- seule l'anomalie de \texttt{PixelSpacing} déjà connue
(\S\ref{sec:bug-ordre-frames}) doit être corrigée par rééchantillonnage, et elle affecte le masque et
l'image de façon cohérente (même résolution des deux côtés, donc un même rééchantillonnage s'applique
aux deux).

\section{Bug critique n°3 (historique) : ordres opposés entre \texttt{volume\_n4.npy} et le nrrd}

Vérifié par corrélation de pixels sur 3 séries : \texttt{volume\_n4.npy[:,:,i]} correspondait à
\texttt{IM\_\{i+1\}} (ordre séquentiel), alors que l'axe des images du labelmap nrrd suit l'ordre
inversé documenté en \S\ref{sec:bug-ordre-frames}. Associer ces deux volumes par le même indice brut
désalignait silencieusement l'image et le label pour toutes les images sauf celle du milieu. Corrigé
dans \texttt{load\_aligned\_pair}. \textbf{Ce bug est aujourd'hui sans objet en pratique} : les
fichiers \texttt{volume\_n4.npy} ont été retirés de \texttt{DATA/} (\S\ref{sec:n4-manquant}), donc
\texttt{load\_aligned\_pair} ne charge plus que du DICOM brut --- mais la correction reste nécessaire
dans le code pour le jour où un volume prétraité y sera réintroduit.

\section{Vérification de l'exclusivité mutuelle}

Vérifiée empiriquement sur les 56 séries : l'exclusivité mutuelle des 6 structures est respectée, à
1 seul pixel de frontière isolé près (dans le fichier à 2 couches déjà cité). Conclusion : une sortie
softmax multi-classe est justifiée, avec une perte combinée Dice + Cross-Entropy.

\section{Composantes connexes et outliers de position}
\label{sec:composantes-connexes}

\subsection{Composantes connexes}

Analyse des composantes connexes par structure : 34 des 336 couples structure$\times$série sont
fragmentés en plus d'un blob (le plus souvent MSC, K1, ASC, PB) --- plausible étant donné les coupes
transversales de ces structures, mais à vérifier visuellement avant utilisation en aval. CLAV et VSC
sont quasiment jamais fragmentées.

\subsection{Outliers de barycentre}

Détection des barycentres 3D aberrants (au-delà de 2$\sigma$ sur un axe) entre patients : \textbf{44
outliers} au total sur les 6 structures $\times$ 3 axes.

\textbf{Hypothèse testée et infirmée (10/08/2026)} : les séries D (droit) et G (gauche) auraient pu
être des images miroir anatomique l'une de l'autre, et leur mélange dans une seule distribution
aurait pu gonfler artificiellement l'écart-type et créer de faux outliers. Le test --- refaire la même
détection séparément par latéralité --- donne \textbf{exactement le même total (44 outliers dans les
deux cas)} : cette hypothèse est rejetée, le signal n'est pas un artefact du mélange D/G.

Les outliers sont en réalité concentrés sur \textbf{3 patients récurrents}, sur les deux latéralités
et plusieurs structures/axes : \texttt{2017-110\_01-0222-V1MR} et
\texttt{2017-110\_01-0223-V1MR/FLASH D} (déjà exclus pour géométrie non résolue, \S précédente) et
\textbf{\texttt{2017-110\_01-0224-V1MR}, qui n'a aucune autre anomalie connue et mérite une
vérification visuelle manuelle sous 3D Slicer avant la Phase 0}.

\todonote{Note statistique : avec 6 structures $\times$ 3 axes $\times$ 2 latéralités = 36 tests
indépendants à 2$\sigma$ sur des échantillons de 25 à 28 séries chacun, quelques faux positifs sont
attendus par pur hasard --- ne pas sur-interpréter les outliers isolés proches du seuil, se concentrer
sur les patients qui reviennent systématiquement.}

\section{Distribution de taille par structure}
\label{sec:taille-structures}

Nombre moyen de voxels annotés par frame et par structure, sur les 51 séries à géométrie résolue :

\begin{center}
\begin{tabular}{lrrrr}
\toprule
Structure & Moyenne & Écart-type & Min & Max \\
\midrule
PB   & 403,3 & 221,9 & 94,2  & 1101,6 \\
K1   & 376,7 & 114,4 & 187,0 & 667,9 \\
CLAV & 368,6 & 119,0 & 174,4 & 632,2 \\
VSC  & 175,9 & 56,3  & 86,2  & 349,0 \\
MSC  & 166,2 & 96,0  & 51,5  & 564,3 \\
ASC  & 94,0  & 27,2  & 37,4  & 150,5 \\
\bottomrule
\end{tabular}
\end{center}

Déséquilibre \textbf{modéré} entre structures : ratio de 4,3$\times$ entre la plus grande (PB) et la
plus petite (ASC). Ça justifie une pondération par structure dans la loss Dice en Phase 1, sans
nécessiter de technique de rééquilibrage agressive (sur-échantillonnage, focal loss, etc.).

\section{Comparaison FLASH vs RSSFP}
\label{sec:flash-vs-rssfp}

Motivation : FLASH (écho de gradient) et RSSFP (état stationnaire) ont des mécanismes de contraste
différents ; la littérature s'attend généralement à ce que RSSFP rende le sang hyperintense (pertinent
pour VSC/ASC).

\textbf{CNR (contraste sur bruit) par structure et par séquence}, sur les 51 séries à géométrie
résolue, calculé de façon méthodologiquement cohérente --- toutes les séries chargées depuis le DICOM
brut via \texttt{load\_aligned\_pair}, sans mélange avec d'anciens volumes N4 précalculés
(\S\ref{sec:n4-manquant}) :

\begin{center}
\begin{tabular}{lrr}
\toprule
Structure & CNR FLASH (moy.) & CNR RSSFP (moy.) \\
\midrule
ASC  & 2,19 & 0,17 \\
VSC  & 2,59 & 0,11 \\
CLAV & 0,24 & 0,11 \\
MSC  & 0,58 & 0,53 \\
K1   & 0,35 & 0,50 \\
PB   & 0,70 & 0,94 \\
\bottomrule
\end{tabular}
\end{center}

\textbf{Avant correction du bug de transposition (\S\ref{sec:bug-transposition})} : la conclusion
provisoire était ``pas de différence systématique forte'' entre FLASH et RSSFP. \textbf{Cette
conclusion est invalidée} et ne doit plus être citée.

\textbf{Après correction}, confirmé visuellement --- sur FLASH, ASC se situe précisément sur une zone
hyperintense de flux nette ; sur RSSFP, la même région est bruitée sans contraste net.

\textbf{L'écart FLASH/RSSFP est-il un effet population robuste, ou un artefact de variance
inter-patient ?} Test : comparer l'écart moyen FLASH$-$RSSFP à l'écart-type inter-patient pour chaque
structure.

\begin{center}
\begin{tabular}{lrrl}
\toprule
Structure & Écart FLASH$-$RSSFP & Écart-type inter-patient & Verdict \\
\midrule
VSC  & +2,31 & 0,43 & écart domine --- effet robuste \\
ASC  & +1,84 & 0,53 & écart domine --- effet robuste \\
CLAV & +0,15 & 0,24 & variance inter-patient domine \\
MSC  & +0,09 & 0,22 & variance inter-patient domine \\
K1   & $-$0,07 & 0,31 & variance inter-patient domine \\
PB   & $-$0,14 & 0,20 & variance inter-patient domine \\
\bottomrule
\end{tabular}
\end{center}

\textbf{Conclusion retenue} : l'avantage de FLASH sur ASC et VSC est un \textbf{effet population
robuste} (il dépasse largement la variabilité entre patients), pas un artefact d'un ou deux cas
particuliers. C'est contraire à l'attente générale de la littérature sur les séquences steady-state et
suffisamment surprenant pour être discuté avec le directeur de thèse (hypothèses : paramètres
d'acquisition spécifiques, agent de contraste sur une seule séquence, spécificités du protocole) avant
d'interpréter la comparaison FLASH vs RSSFP au niveau modèle en Phase 1. Pour CLAV, MSC, K1 et PB, la
variance inter-patient domine l'écart entre séquences : \textbf{aucune conclusion FLASH vs RSSFP ne
peut être tirée de l'EDA seule pour ces 4 structures} ; la vraie comparaison vient des métriques des
deux modèles spécialisés entraînés en Phase 1 (FLASH-only et RSSFP-only, pas un choix d'une séquence
de référence unique --- voir chapitre Protocole expérimental).

\section{État des volumes N4}
\label{sec:n4-manquant}

Seuls 32 des 56 dossiers patient/série disposaient de \texttt{volume\_n4.npy} (vérifié le 08/08/2026).
Décision prise dès l'EDA : ne pas construire un pipeline mélangeant fichier précalculé (quand présent)
et DICOM brut (sinon), car ça biaiserait le prétraitement de façon inconsistante entre patients ---
dangereux pour une analyse d'interprétabilité.

\textbf{Décision appliquée le 10/08/2026, pas seulement documentée} : les 32 fichiers
\texttt{volume\_n4.npy} existants ont été déplacés hors de \texttt{DATA/} vers
\texttt{\_archive\_volume\_n4/}. \texttt{load\_aligned\_pair} retombe donc systématiquement sur le
DICOM brut pour les 56 séries, ce qui a aussi rendu la comparaison FLASH vs RSSFP
(\S\ref{sec:flash-vs-rssfp}) méthodologiquement cohérente --- elle ne mélange plus des séries
N4-corrigées et des séries brutes, contrairement à une version antérieure de cette EDA. La correction
N4 elle-même (\texttt{src/preprocessing.n4\_bias\_correction}, SimpleITK, shrink-then-apply-full-res)
est écrite mais reste à brancher uniformément dans le pipeline de Phase 0.

\section{Table de synthèse}

Le notebook se conclut par une table de synthèse unifiée par série, croisant inventaire, statut de
géométrie, métadonnées, taux de fragmentation et CNR moyen. \textbf{51 séries sur 56 sont utilisables
automatiquement} pour la Phase 0/1 ; les 5 restantes (patient \texttt{0222} entier +
\texttt{0223/FLASH D}) nécessitent une vérification manuelle sous 3D Slicer.
